% =====================================================================
%  R. Nielsen, DEN STÆRKERES RET.  Tale ved Reformationsfesten 1863 for Universitetet.
%  Kjøbenhavn: Gyldendal (F. Hegel); Thieles Bogtrykkeri, 1863.  15 pp.  Fraktur.
%  The speech given at the University of Copenhagen's Reformation festival, 1863.
%
%  TRANSCRIPTION -- GENERATED; DO NOT EDIT BY HAND.  Rebuild, chained (&&), from the repo root:
%    python3 ebook-method/staerkeres-ret/draft.py --emit --force
%      this template + the scan's text layer (Google's OCR in the Google Books scan of the
%      British Library copy 3911.aa.46., books.google.com/books?id=QDZdAAAAcAAJ;
%      ~/bibliotek/Nielsen, Rasmus/1863-den_Stærkeres_Ret_Tale_ved_Reformations.pdf, see SCANS.tsv);
%    python3 ebook-method/collate/check.py staerkeres-ret --witness ocr --apply
%      the decisions in decisions.tsv, taken against tesseract's independent reading of the images;
%    python3 ebook-method/staerkeres-ret/oepass.py --apply
%      word-fixes.tsv (the page check's word corrections, each seen in a crop);
%    python3 ebook-method/staerkeres-ret/finish.py --apply
%      the drop capital, a dash, the printed rule on p. 14, PRINTED AS IS marks (each read at the image).
%
%  WORD-ACCURACY (2026-10-08).  No e-book and no earlier transcription exists.  The draft (Google's OCR layer) was
%  collated word by word, with punctuation and word division, against tesseract (tessdata_best script/Fraktur), with
%  deu_latf and frak2021 as third and fourth readings: 284 disagreements.  70 settled by check.py's own rules, 161 by
%  rules.py (each decision named; a blind audit of 22 of them, mixed into the crops, found one rule error, mended:
%  „samler“ p. 15), 50 at the image in BLIND crops read by a Sonnet reader who saw neither OCR reading (2 of 2 planted
%  controls read right), 3 by hand from the reader's notes.
%
%  FULL PAGE CHECK (2026-10-08, pagecheck.py with every page in the sample: the title page and pp. 3-15, 3857 words).
%  Every page was transcribed BLIND at the image by a second Sonnet reader and compared word by word; every difference
%  (46) was settled in a crop read in the main session.  Transcription errors found and mended: 16 of words (4.2 per
%  1000) and 1 of punctuation; the other 29 differences were the reader's slips or notation.  Of the 16, 8 belong to the
%  two classes of fault that the layer AND tesseract share (see skjaebne-og-forsyn): the Fraktur Ø read as D (Øieblik
%  x5, Øiemed) and gj read as gi (gjøre x2, where the rule frak had overruled the layer); 4 more are ø read as o (Møde,
%  høiere x2, Høihed), which every OCR misses; then Hær (layer: Har), Glosser (Glosfer), Illusionen (JIllusionen), and
%  the drop capital P read twice.  The dash p. 9 („Hvad er — paa“) had been dropped by the layer.  Since every page was
%  read, these are all the errors the check could find; what it cannot see is a reading in which the transcription and
%  the blind reader agree on the same fault.
%  Paragraphing: 29 indented paragraph starts, as the page reader counted (draft INDENT_MAX measured to 5.0).  Emphasis:
%  the page reader saw no letter-spacing or bold in the text (the title page only).
%  PRINTED AS IS: „Ostandelse“ for Opstandelse (p. 9); „Oltids“ for Oldtids (p. 15).  „disse dristige, driftige
%  Forventninger“ (p. 13): the second word has f, not long s (600 dpi crop) -- not a doubled word.
%
%  PAGE MAP.  printed p = PDF page - 6 (p. 3 = PDF 9, p. 15 = PDF 21); folios 4-15 printed, p. 3 has none (a
%  handwritten „64“ at its foot).  PDF 7: title, p. [1]; PDF 8: p. [2], blank; PDF 22: the British Library's date stamp.
%  No signatures, no footnotes.
%
% =====================================================================
\documentclass[12pt,a4paper,openany]{book}
\usepackage[utf8]{inputenc}
\usepackage[T1]{fontenc}
\usepackage{libertinus}
\usepackage{graphicx}
\DeclareUnicodeCharacter{0254}{\reflectbox{c}}  % ɔ (U+0254) = reversed c
\DeclareUnicodeCharacter{A75B}{r}   % r rotunda (ꝛc. = etc.)
\usepackage[danish]{babel}
\usepackage[protrusion=true,expansion=true]{microtype}
\usepackage{setspace}
\usepackage[margin=1.3in]{geometry}
\usepackage{fancyhdr}
\usepackage{hyperref}
\hypersetup{pdftitle={Den Stærkeres Ret}, pdfauthor={Rasmus Nielsen}, hidelinks}
\setlength{\headheight}{15pt}
\pagestyle{fancy}
\fancyhf{}
\fancyhead[LE,RO]{\thepage}
\fancyhead[C]{\textit{R. Nielsen: Den Stærkeres Ret (1863)}}
\renewcommand{\thefootnote}{\ifcase\value{footnote}\or *)\or **)\or ***)\else
  \arabic{footnote})\fi}
\setstretch{1.3}
\usepackage{marginnote}
\newcommand{\opage}[1]{\marginnote{\footnotesize\textit{[#1]}}}
\newcommand{\apage}[1]{\marginnote{\footnotesize\textit{[#1]}}}
% headings as printed (read at the image, .parts/notes/r-struct.tsv): the Indledning and the part titles in large
% Fraktur, the part letters A./B./C. and section numerals I.-IV. in bold antiqua, section titles in bold Fraktur; a short
% rule under each part title; in B, sections I and II take a bold Scripture verse as their title (\motto)
\newcommand{\overskrift}[1]{\begin{center}{\Large #1}\\[0.3em]\rule{3em}{0.4pt}\end{center}\addcontentsline{toc}{section}{#1}}
\newcommand{\afdeling}[3][]{\begin{center}\ifx&#1&\else#1\\[0.6em]\fi\textbf{#2}\\[0.3em]{\Large #3}\\[0.3em]\rule{3em}{0.4pt}\end{center}\addcontentsline{toc}{section}{#2 #3}}
\newcommand{\delafsnit}[1]{\begin{center}{\Large\textbf{#1}}\end{center}\addcontentsline{toc}{subsection}{#1}}
\newcommand{\stykke}[2]{\begin{center}\ifx&#1&\textbf{#2}\else\textbf{#1}\ifx&#2&\else\\[0.2em]\textbf{#2}\fi\fi\end{center}}
\newcommand{\motto}[1]{\begin{quote}\textbf{#1}\end{quote}}

\begin{document}

% --- PDF 7, the title page, p. [1], read at the image 2026-10-08 (page reader + main session): Fraktur throughout
% but the year (antiqua); „Tale“, „Universitetet“ and „Thieles Bogtrykkeri.“ letter-spaced.  Not printed, so not
% transcribed: a handwritten „1624“ at the head, a pencilled „K“ and underline, the British Museum's stamp.
% The verso, p. [2] (PDF 8), is blank.
\begin{titlepage}
\begin{center}
\vspace*{4em}
{\Huge\textbf{Den Stærkeres Ret.}}\\[2em]
{\Large\textls[200]{Tale}}\\[1em]
{\large ved Reformationsfesten 1863}\\[0.8em]
{\small for}\\[0.8em]
{\large\textls[200]{Universitetet}}\\[0.8em]
{\small af}\\[0.8em]
{\large\textbf{R. Nielsen.}}\\[8em]
{\small\textbf{Kjøbenhavn.}}\\[0.5em]
{\small Forlagt af den Gyldendalske Boghandling (F. Hegel).}\\[0.5em]
{\footnotesize\textls[200]{Thieles Bogtrykkeri.}}\\[0.5em]
\textit{1863.}
\end{center}
\end{titlepage}

% ===== TEXT (pp. 3--15) =====
% --- p. 3 ---
\setcounter{footnote}{0}%
\opage{3}Paa Reformationens Vei, dens Strid, dens Gang, dens
Kamp og Seir, lyser en Tanke, --- den gamle Tanke om
Kraft og Svaghed, om Ret og Magt; og denne Tanke er
en Dom, en Verdensdom om Ret og Magt: den Stærkere
vil have Ret, den Svagere maa give efter.

Ja, det er en Verdensdom: den Stærkere vil have Ret.
„Vi ville have Ret, thi vi have Magt!“ saa lød det i Trusler
fra Rom; „vi skulle vel faae Magt, saa sandt vi have Ret!“
saa lød det i Svar fra Wittenberg. Saa høres fra alle
Sider de Stridendes Raab paa Ret, paa Magt og paa
Forlig mellem Ret og Magt; det er som vilde de alle det
Samme. Og dog er her en Forskjel, en væsentlig Forskjel.
Thi det gjør jo dog en væsentlig Forskjel, om vi forstaae
Rettens Eenhed med Magten saaledes, at det er den sandselige,
haandgribelige Magt, der uden alle Omsvøb er Ret, eller
saaledes, at det er den hellige, den uforkrænkelige Ret, der i
sit Væsen evig er Magt; det giver en Verdens Forskjel, om
vi bedømme Styrken saaledes, at den Stærkere har Ret,
alene fordi han er stærk, eller saaledes, at han i Sandhed
kun er den Stærkere, forsaavidt han har Ret.

Her er Forskjel paa Ret og her er Forskjel paa Magt:
her er physisk Magt, og her er aandelig Magt; men hvilken
af disse Magter man end vil forvandle til Ret, dersom det
ikke betyder, at man med det Samme vil forvandle Ret til
Magt, saa bliver dog Magten misbrugt.

% --- p. 4 ---
\setcounter{footnote}{0}%
\opage{4}Der er Elendighed i Verden, der er Jammer og Nød,
naar den physiske Magt bliver misbrugt; thi det er tungt,
ubeskrivelig tungt, at bøie sig for den raa Magt; det er tungt,
ubeskrivelig tungt, at erkjende den Stærkeres Ret, naar den
Stærkere kun er en Røver; det er tungt, ubeskrivelig tungt,
at bære Undertrykkelsens Aag, og tungt at opveie Frihed
med Guld. Og Frihed veiet mod Guld! hvad er vel Guldets
Vægt, naar en Brennus kaster sit Sværd paa Vægtskaalen?
væ victis! --- ja, vee de Overvundne.

Og dog er Elendigheden større endnu, naar den aandelige
Magt bliver misbrugt; naar Bedrag omskygger det
Hellige, naar Sjælene fanges og Tankerne bindes og Sandhed
svinder og giver et Skin som Solens matte Skin i en
Formørkelse. Vel er Krigen med sine Rædsler forfærdelig;
men Pesten er dog endnu forfærdeligere; og vel er Trældom
i det Udvortes en sørgelig Ting, men Trældom i det Indvortes,
aandelig Trældom, Samvittighedens og Tankernes Trældom,
er Pesten selv, den skumle Pest: denne sorte Død er Fortabelse.

Vi tør nu ikke sige, at den katholske Kirke begyndte med
at bøie Retten og misbruge Magten; thi, timelig seet, begyndte
den ikke med Magten. Den begyndte ikke i Pragt
og Herlighed, men i Forsagelse, i Selvfornegtelse og udvortes
Ringhed. Det begyndte ikke med, at en Christi
Statholder og Fyrsternes Formynder sad høit paa St.
Peters Stol og slyngede Banstraalen; nei, Begyndelsen var
i Sagtmodighed, i Opoffrelse, i Taalmodighed. Fredens
Sendebud gik ud i Verden, ud til Folkene i den vilde
Strid, at bringe dem Forsonerens Hilsen og forlige Fienderne;
saaledes var Begyndelsen.

Hvad Forholdet er mellem Magt og Ret, kan sees af
denne Begyndelse. Havde hine Sendebud ikke været forvissede
om, at der gives en uforkrænkelig Ret, der er stærkere
end den raa Magt: de havde ikke bragt Forsonerens Bud%
% --- p. 5 ---
\setcounter{footnote}{0}%
\opage{5}skab, de havde ikke vovet sig ud blandt de vilde Horder. De
kom som Værgeløse, men dog saa vel udrustede; de kom som
Svage, og dog saa stærke! thi den, der kan herske over sit
Sind, er stærk, stærkere end den, der kun kan indtage Staden;
og den, der kan bede for sine Fiender, er stærk, stærkere end
den, der kun kan myrde sine Fiender, og den, der kan offre
sit Liv for Sandhed og Ret, er stærk, stærkere end den, der
kan offre Folkeslag og Riger for indbildt Hæder og Magt.

Saa gik da hine Sendebud, Enhver ad sine Veie, ud
i Verden. Der kommer en Fremmed til Krigernes Leir;
det er en Vaabenløs, der vil stride med de Bevæbnede, en
Svag, der dog vil overvinde de Stærke: hvo skal i denne
Strid give efter? Det bliver ikke den Vaabenløse, der giver
efter for de Bevæbnede; nei, det bliver de Bevæbnede, som
denne Gang maae give efter for den Vaabenløse. Det er jo
ikke den Fremmede, der skjælver for Herskerens Blik, bliver
angst for det blinkende Sværd og kaster sig ned for Kongens
Fødder og beder om Beskærmelse og Naade; nei, han skjælver
ikke for Herskerens Blik; han forkynder sit Budskab; han er
slet ikke angst for det blinkende Sværd, thi han forkynder sit
Budskab; han beder ikke om Naade, han tilbyder Naade.
Nu forstumme de vilde Raab; det bliver stille i Leiren:
Krigeren hører den Fremmedes Budskab, Kongen, den ellers
Ubøielige, bøier sit Hoved, lytter og grunder og føler sig
underlig tilmode. Saa er det da ikke Hæren, denne hele
Hær af stærke Mænd, fuldblodige Krigere, som i dette Møde
har Magten, den høiere Magt; nei, den Fremmede med sit
Budskab, Ordets Mand med Blikkets Ild, den blege Munk
med Crucifixet, ham er det, som i dette Møde har Magten,
den høiere Magt; ham er det, som i dette Møde bærer
Magtens Tegn, det Magtens Tegn, som evigt er Retfærdighedens
Tegn. See dette Magtens Tegn! det gjør den Svage
stærk, den Stærke svag; thi Kongen knæler med sin Hær for
% --- p. 6 ---
\setcounter{footnote}{0}%
\opage{6}Magtens og Retfærdighedens Tegn. Saa er det da en
Dom, en Verdensdom om Ret og Magt: den Stærkere vil
have Ret; den Svagere maa give efter.

Der er Forskjel paa Ret og der er Forskjel paa Magt;
der vil i dette Timelige altid dannes og digtes, spidsfindig
digtes en Ret, som dog i evig Forstand er Uret; der vil i
det lavere Virkelige altid røre sig en Magt, en opblæst, raa,
paatrængende Magt, som dog i evig Forstand er Afmagt.
Skal nu det Timelige ikke destomindre tvinges til at bøie sig
under det Evige; skal Retfærdigheden selv, den hellige Tanke,
som evig er Magt, gribe ind i Tidernes Hjul og bevæge
Virkeligheden: da maa dens Aabenbaring være mere end et
stort Syn i et stort Øieblik; da maa dens Gjennembrud
være Mere end et Blus i Skyen, et Lyn i Natten eller et
flagrende Lys i Stormen; det Guddommelige vil da ikke
indskrænke sig til at overraske Folkene ved Tegn fra Himmelen,
men nedstige til Menneskene, nedlade sig til menneskelige
Formaal, menneskelige Vilkaar, og tage Bolig blandt
Menneskene.

Men naar det Guddommelige saaledes fra sin Side
rører sig, for at rense, styrke og forklare det Menneskelige,
kan det heller ikke feile, at jo et Menneskeligt fra sin Side
ogsaa vil røre sig, for at svække, smitte og formørke det
Guddommelige.

Folkelig praktisk drog den katholske Kirke med dristig
Haand det Oversandselige ned i Sandsernes Verden; folkelig
praktisk lod Catholicismen det Usynlige iføre sig synlig
Skikkelse og gav det Allerhelligste legemlig, ja haandgribelig
Nærværelse. Det var et Skridt nedad; men dog i menneskelig
Forstand tillige et Skridt fremad. Paa det sandselige phantasirige
Naturmenneske virker Aanden ved synlige, sandselige
Midler. Middelalderens Folk med de ideale Drømme, de
dybe Længsler og stærke Lidenskaber trængte til aandelige
% --- p. 7 ---
\setcounter{footnote}{0}%
\opage{7}Formyndere, til Tugtemestere med guddommelig Fuldmagt;
Næverettens Tid, der fordrede Guds-Domme, forlangte at
høre --- Guds egen Røst, forlangte at see --- Guds Statholder
selv nærværende paa Jorden. Det var en egen Blanding af
Barbarie og Fromhed, af Trods og Ydmyghed, raa Sandselighed og asketisk Selvfornegtelse i utallige Afskygninger af
Vildfarelser og Smerter, Forbrydelser og Anger og Taarer,
der i hiin eventyrlige Tidsalder gjorde Theokratiet nødvendigt.
Med Theokratiet fuldendtes Hierarchiet. Saa reiste sig den
høie Kirke; Middelalderens Kirke paa Oldtidens Gruus,
paa Hedenskabets mægtige Ruiner; Romerkirken med de
dunkle Hvælvinger, med den dybe Choral og det mystisk
alvorlige Sprog; Pavekirken med de Alt formaaende Naadegaver,
et uberegneligt Overskud af gode Gjerninger, en
uvurdeerlig Skat af skriinlagte Helgenes Been og af alle de
Helliges Fortjenester. Saaledes reiste sig den høie Kirke:
det er Middelalderens Aand i Middelalderens Stiil, denne
symbolske Ritus, denne billedrige, sindbilledlige Cultus med
Skriftemaal og Litanier, med Alterlys og Virakskyer, Præster
og Sakramenter.

Vistnok var her i alt dette en Forvanskning af det Oprindelige,
en stedse voxende Sandhedsformørkelse. Evangeliets
Lys stod allerede længst i Skygge bag Kirkens Piller; Menighedens
Herre, Forbilledet selv, var glemt af en Skare Tilbedere,
der helst vilde knæle for Himmelens Dronning. Men
trods denne Sandhedens Formørkelse var Kirkens høiere
Magt dog endnu altid begrundet paa høiere Ret: det Menneskelige
var opløftet ved det Guddommelige; vidunderlige Kræfter
rørte sig i Sind og Tanker. Kirken i sin Magt og Vælde
havde allerede skabt sig en ny Himmel; der indtraadte nu et
Vendepunkt, hvor det ret skulde vise sig, at den ogsaa formaaede at skabe en ny Jord; dette Vendepunkt er Korstogenes
Begyndelse.

% --- p. 8 ---
\setcounter{footnote}{0}%
\opage{8}At Europa i det 11te Aarhundrede trængte til en Opvækkelse,
er umiskjendeligt. Under mangehaande Trængsler
og mørke Anelser arbeidede det med tunge, uklare Tanker paa
at finde en Vei gjennem Virkeligheden, et Bedrifternes Maal,
en Mulighed til at oprede Timelighedens Virvar. Det var
en qvalfuld Tilstand. Vel var der en Salighed i Himlen
for dem, der undgik Fordømmelsen; men paa Jorden var
der intet Godt, thi Jorden var en Jammerdal. Saa gik
der da en Angst igjennem alle Sjæle; der skete Tegn i Sol
og Maane; Folkene vansmægtede i Rædsel og Forventning, thi
Himmelens Kræfter rørte sig: Dommedag var nær forhaanden.

Skulde dette Jordiske --- endnu engang --- faae virkelig
Betydning, da kunde det kun skee derved, at det umiddelbart
kom til at tjene det Himmelske; skulde der gives et
høiere Maal for Bestræbelser i det Timelige, da maatte der
være en Vei paa Jorden, en Hæderens, Handlingernes og
Bedrifternes Vei, der førte lige til Himmelen. „Hvad skal
jeg gjøre, at jeg kan arve det evige Liv?“ Det var ikke en
enkelt riig Yngling, som i en Slags Overilelse spurgte saaledes;
thi tusinde og atter tusinde Ynglinge, rige paa Daadskraft
og ridderligt Mod, grundede hver for sig paa det
Spørgsmaal: „hvad skal jeg gjøre, at jeg kan vinde et  % PRINTED AS IS: this quotation is not closed (after „Navn?“; page reader and layer agree)
evigt Liv, et udødeligt Navn? Det var ikke en lille Flok af
Ørkesløse, som, efterdi de i Grunden slet Intet vilde gjøre,
fik Tiden til at gaae med det Spørgsmaal: „hvad skulle vi
gjøre, at vi kunne arbeide Guds Gjerning?“ det var Europas
Folk, som i Trængsler og Nød, i Gjæring og Spænding,
ikke vidste, hvad de skulde gjøre.

Men Kirken vidste det; den høiere Magt er atter her
en høiere Ret; thi Paven vidste, hvad de skulde gjøre. Paa
Kirkeforsamlingen i Clairmont blev det forklaret for de lyttende
Skarer, hvad Christenheden skulde gjøre. „Befrie
Jerusalem!“ saa lød Forklaringen fra Pavens Mund; befrie
% --- p. 9 ---
\setcounter{footnote}{0}%
\opage{9}Jerusalem og myrde Korsets Fiender, det var --- til Syndernes
Forladelse, til Sjælefred og Salighed --- den store
Ridderdaad, den Gjerning, som Guds Folk da skulde gjøre.
„Befrie Jerusalem!“ saa lød det over Vestens Lande; nu
var da Veien funden, Maalet sat, det høie Maal. „Befrie
Jerusalem!“ nu var det Alle klart, at Pavens Tanke var
Guds Tanke og Pavens Villie Guds Villie, ja, saa maatte
Pavens Bud da ogsaa være det Magtens Bud, Hærskarerne
adløde. „Befrie Jerusalem!“ det var Opvækkelse til Liv,
Ostandelse endnu før Dommens Dag.  % PRINTED AS IS: „Ostandelse“ for Opstandelse (p. 9; crop J13, all three OCRs agree)

Men Tiden gaaer, Erfaringen kommer. Det første Syn
af Østerlandene, af Palæstina, af Jerusalem, hvor var dog
dette første fjerne Syn i Phantasiens Farvepragt opløftende
og skjønt og herligt i Sammenligning med al den haarde,
knudrede, kummerlige Virkelighed, der siden Skridt for Skridt
paatrængte sig de endnu Uerfarne. Jo mere overspændt en
Forventning er, desto større er ogsaa den Skuffelse, der
følger efter. Den lange Strid om Opstandelsens Grav er
et romantisk Epos, et eventyrligt Ridderspil med phantastiske
Afvexlinger af Seire og Nederlag, glimrende Forhaabninger
og bittre Skuffelser, opblussende Begeistring og trykkende
Mismod. Det varede i to hundrede Aar, men saa var
Illusionernes Tid forbi.

Men er Illusionernes Tid forbi, saa er det væsentlig
forbi med Middelalderen, thi Middelalderen er Illusionernes Tid. Hvad er ---
paa Historiens Vei --- vel Illusionen? Ideernes Fata Morgana, en Phantasiens Luftspeiling,
fortryllende, naar den i samme Syn, i eet og
samme Billede kan sammensmelte Sandhed med Bedrag.
Saalænge Sandhedselementet endnu er kraftigt til at tvinge
Bedraget, er Illusionen som en Aabenbaringsform for ideale
Magter og kan i denne Charakteer just tjene Begeistringen;
men efterhaanden som Forholdet vender sig om, alt som
Bedraget faaer Overhaand og fortærer Sandhedselementet,
% --- p. 10 ---
\setcounter{footnote}{0}%
\opage{10}vil Idealet vige bort og efterlade sit meduseagtige Modbillede,
et ret afskrækkende Vrængebillede.

Just som det Tidspunkt var indtraadt, da Pavedømmet
havde udtømt sine aandelige Kræfter og dermed tabt sin
væsentlige Berettigelse, kunde Paven med de tvende Sværd
jo vende sig mod Syd og Nord og Øst og Vest, erklære sig
for Verdens Herre, en myndig Hersker over Folk og Stater;
hans Ret, hans kanoniske Ret med Glosser og Extravaganter
blev da fuldstændig bragt i System. Det er et storstilet Rethaveri
her fortsættes; en oprørende Strid mellem Magt og
Ret vil nu for Alvor tage sin Begyndelse.

En ikkun verdslig Magt maa dog altid, dersom den da
i nogen Maade skal bygge paa en Retstilstand og give sig
endog blot Skin af Ret, finde sig i visse Skranker, binde
sig til visse ubrødelige Vilkaar; men med en geistlig Magt,
en hellig, ufeilbar, guddommelig Magt, en Magt, der længst
har krænket den evige Sandhed og gjort det Allerhelligste til
Middel for lave verdslige Øiemed, er det en anden Sag:
her bliver i Ordets sørgelige Mening den høieste Uret
høieste Ret.

For at forestille sig Kirkens geistlig-verdslige Magt i
dens overordentlige Storhed og overordentlige Misbrug, maa
man ret betragte den næsten uoverskuelige Mængde Anmasselser
og Vilkaarligheder, som denne Magt, denne alene ufeilbare
Magt, har været istand til at byde Nationer og Fyrster;
det er en aldrig afbrudt Række af forargelige Optrin, hvormed
dens hellige Herskesyge, den hellige Gjerrighed, det hellige
Rænkespil og den hellige Frækhed har formaaet at udfylde
det mærkværdige Tidsrum, der ligger imellem Korstogenes
Slutning og Reformationens Begyndelse.

Den Stærkere vil have Ret; men her er Forskjel paa
Ret og Forskjel paa Styrke. Saa stærk var Romerkirken
midt i sin Fordærvelse, at hverken de hæftige Storme, der
% --- p. 11 ---
\setcounter{footnote}{0}%
\opage{11}under det pavelige Schisma rasede i dens Indre, eller de
mange Grundstød, som de humane Magter: Staten, Videnskaben,
Nationernes opvaagnende Selvfølelse og den fremadskridende
Cultur, i stedse stigende Reactioner bibragte den
udenfra, vare istand til at kue den Vældige. Eet, men ikkun
Eet, var endnu tilbage, der kunde tvinge den ellers Ubetvingelige; Eet, men ikkun Eet, formaaede at gjennembryde
hiint uhyre Væv af Underfundighed, Løgn og Bedrag; dette
Ene var Sandheden i al Eenfoldighed, den længe forhaanede
Sandhed, Retfærdighedstanken, den ideale Ret, der i sit
Væsen evig er Magt.

Men naar her virkelig skal begyndes med den ideale
Ret, altsaa oprindelig og fra Roden af, da kan her ikke begyndes
med haandgribelig Overlegenhed eller med forud beregnet
Sandsynlighed for en afgjørende Seir; skal her strides
for Sandheden, ubetinget for Sandhedens evige Ret, da
maa her atter begyndes i Svaghed, i timelig Ubetydelighed,
i udvortes Ringhed: saaledes begynder Reformationen.

„Vi ville have Ret, thi vi have Magt,“ saa lød det i
Trusler fra Rom; „vi skulle vel faae Magt, saa sandt vi
have Ret,“ saa lød det i Svar fra Wittenberg. At den hellige
Ret i sit Væsen evig er Magt, var just den Tanke, der besjælede
Luther; det var en Samvittighedstanke, der hævede
ham høit over Menneskefrygt, en Guddomstanke, der gav
ham Kraft og Mod, en saadan Kraft, et saadant Mod, at
han i hvert Afgjørelsens Øieblik --- og med hvor mange Afgjørelsens
Øieblikke er ikke hans Levnetsløb betegnet?
--- altid fandt det afgjørende Ord og stadfæstede Ordet ved afgjørende
Handling. Paa Rigsdagen i Worms forhandles da
et stort Retsspørgsmaal; paa Rigsdagen i Worms ere Magthaverne
forsamlede. Betragter engang Forsamlingen, det er
en talrig, en glimrende Forsamling af Riddersmænd, Prælater,
Fyrster og Churfyrster med Keiseren selv, den nylig Kronede,
% --- p. 12 ---
\setcounter{footnote}{0}%
\opage{12}i sin Midte: Alt hvad det hellige romerske Rige har at opvise
af Glands og Hæder, af Pragt og Herlighed, er her tilstede;
og betragter saa den ringe Munk, den haardt Anklagede,
den to Gange Bandlyste, der nu træder ind i Forsamlingen
og skal føre sit Forsvar: hvilken Modsætning i det
Udvortes af Ringhed og Høihed, af Afmagt og Magt maa
ikke paatrænge sig os ved dette Syn; og saa --- hvilken
Modsætning i det Indvortes af Høihed og Ringhed, af Magt
og Afmagt vender ikke Forholdet om og forvandler os Indtrykket
af dette Syn!

Det er da igjen den udvortes, haandgribelige Magt, der
truer med at knuse den Værgeløse, og det er igjen den
Værgeløse, der vil stride med de Bevæbnede, en Svag, der
dog vil overvinde de Stærke. Saa vilde Luther bryde Løgnens
Aag; saa lød hans sidste Svar høit i den lyttende Forsamling:
„Eftersom da Eders Keiserlige Majestæt, Churfyrstelige
og Fyrstelige Naader begjære et eenfoldigt, slet og
ret Svar, saa vil jeg give et, som hverken har Horn eller
Tænder, nemlig saaledes: Med mindre jeg ved den hellige
Skrifts Vidnesbyrd eller med aabenbare, klare og indlysende
Grunde bliver overvunden og overbeviist (thi jeg troer hverken
Paven eller Concilierne, da det ligger aabenbart for Dagen,
at de ofte have feilet og modsagt hverandre), saa at jeg ved
selve de Sprog, som jeg har anført og brugt, bliver overtydet, og min Samvittighed fangen i Guds Ord, hverken kan
eller vil jeg kalde Noget tilbage, thi det er hverken sikkert
eller raadeligt at gjøre Noget imod sin Samvittighed. Her
staaer jeg, jeg kan ikke Andet; Gud hjælpe mig. Amen!“

Det var i Afgjørelsens Time et afgjørende Ord, og
Ordet var en afgjørende Handling. Hvo kan nu tvivle om,
hvad Luther vil, og hvad det er, hvorfor han strider? For
Sandheds hellige Skrift, for de aabenbare, klare og indlysende
Grunde og for Samvittighedens uforkrænkelige Ret!
% --- p. 13 ---
\setcounter{footnote}{0}%
\opage{13}Vel staaer han der alene i de Mægtiges Forsamling; men
dog ikke ene som den, der er forladt af Alle. Vistnok har
han ikke forud sikkret sig en physisk Understøttelse, en Hær af
Forbundstropper, der kunde rykke frem som fra et Baghold
og komme ham paa givet Vink til Undsætning; men derfor
er han ikke uden Hjælp og Bistand. I Farens Øieblik, just
da han handler som den, der ikke mere kan underhandle, i
Farens Øieblik bliver der i det Forborgne arbeidet for hans
Sag, paa tusinde Steder, i tusinde Maader virket for ham,
underhandlet til hans Hjælp og Bistand. Thi hvormange
nye Tanker, hvormange hidtil dunkle, undertrykte Forestillinger, selv i de kronede Hoveder, hvilke ubeskrivelige
Anelser og dristige Forventninger blive her ikke vakte, kaldte
frem som ved et Trylleslag af Luthers stærke Ord. Og disse
nye Tanker, disse frisk udsprungne Forestillinger, disse dristige,
driftige Forventninger, de strømme ham imøde, de tale hans
Sag, de stemme for ham, de samle sig i Kreds om ham,
en mægtig, en usynlig Hær, en Forbundshær af klare Grunde,
stærke Sympathier. Vel har nu Løgnen ogsaa sine Sympathier;
men Sandhed, Ret og klare Grunde --- hvordan
end Verden dreier sig --- de ville altid virke for den gode
Sag; de ville samle sig i Farens Stund, en mægtig Hær
af skjulte Kræfter, stærke Sympathier.

Vistnok var det en ulige Kamp --- af Lysets Straaler
imod Mulm og Mørke, eenfoldig Sandhed mod et Væv af
Løgn; det var ulige Kamp --- med Aandens Vaaben, Ordets
Kraft alene imod physisk Magt og Rænkespil og snedig
Fanatisme; det var ulige Kamp, men i ulige Kamp skal
Styrken prøves. Og den blev prøvet Skridt for Skridt og
Punkt for Punkt, og Prøvelsen var streng; men aldrig
svigtede dog Sandheds hellige Skrift, og aldrig svigtede
de klare Grunde med den stærke Ret. Ja, den blev prøvet,
Prøvelsen var tung; men hiint Retfærdighedens Tegn, det
% --- p. 14 ---
\setcounter{footnote}{0}%
\opage{14}Tegn, der gjør den Svage stærk, den Stærke svag, var atter
Magtens Tegn paa Farens Dag, et Seierstegn i Kampens
Storm, da Luthers Røst lød høit i Stormen: „Hører Du,
Pave, ikke du allerhelligste, nei, du allersyndigste, at Gud
fra sin Himmel skal snarligen styrte din Stol og sænke den
i Helvedes Afgrund!“

Det er ikke Formastelse, det er ikke Overmod og det er
ikke blinde Lidenskaber, der tale til os igjennem disse Trusels
Ord; nei, det er Luthers Tro, det er hans levende og klippefaste
Overbeviisning, at den hellige Ret i sit Væsen evig er
Magt. Det er da heller ikke Paven i Rom alene eller
Pavedømmet alene, hvorpaa Truslen sigter; thi den sigter
paa enhver oprørende Anmasselse; den sigter paa enhver
vilkaarlig og despotisk Magt, som ved et Overmaal af Løgn
og Rænker indbilder sig at kunne formørke Alt og trodse
Alt, baade Sandheds hellige Skrift og de klare Grunde og
den oprindelige stærke Ret: over enhver saa blind, saa uretfærdig
Magt i Tiden har Martin Luther lyst Forbandelse.

\begin{center}\rule{3em}{0.4pt}\end{center}

Den Stærkere vil have Ret: det er en Dom, en Verdensdom!
Men hvad er en Verdens Dom uden en Verdens
Dommer? Hvad er vel al vor Strid om Ret og Magt i
dette Timelige? en vaagen Drøm, et Skyggespil, dersom der
ikke er en evig Ret, en hellig Magt, der kjender Alt og
styrer Alt og aabenbarer sig i Alt som Verdens Dommer.

I din Varetægt, Alt seende Forsyn, almægtige Verdens
Dommer, befale vi da vort elskede Fædreland med Landsfaderen,
vor folkelig trofaste, dyrebare Konge, Kong Frederik
den Syvende, og hele det kongelige Huus. Din Aand, som
er Viisdoms og Forstands Aand, Raads og Krafts Aand,
være over Fædrelandets troe Mænd i Kongens og i Folkets
Raad, at de Institutioner, som værne om Friheden og det
borgerlige Livs rige, alsidige Udvikling i aandelig saavel som
% --- p. 15 ---
\setcounter{footnote}{0}%
\opage{15}i timelig Retning, maa bevares, omfredes og benyttes til i
stedse videre Omfang at forlige det Stridige og forene det
Adskilte. Og naar vi derhos haabe og bede, at denne Høiskole
fremdeles maa blomstre og aldrig forlades af den
levende Sands for Sandhedens Undersøgelse og Sandhedserkjendelse,
der alene knytter varige Baand mellem Lærere og
Tilhørere: da tænke vi herved tillige paa enhver høiere
Stræben i Kirke og Skole, i Videnskab og Kunst, at alle
Aandens Kræfter tilsidst maae rettes imod samme Maal; vi
tænke paa ethvert redeligt Arbeide, enhver god og hæderlig Gjerning
her i Landet, at vi Alle maae have Fredens rige Velsignelse,
saalænge det kan forundes os at leve i Fred.

Men er Fred umulig, slaaer Afgjørelsens Time,
kommer Kaldelsen til os, da vi i Gjerning skulle vise, at vi
endnu --- endnu i den ellevte Time --- have bevaret Luthers Tro,
hans klippefaste Overbeviisning, at der er en uforkrænkelig Ret,
som i sit Væsen evig er Magt: da være Du Almægtige, i
den ulige Kamp, vort Skjold og Værge og vor faste Borg,
som Du var Luthers faste Borg, hans Skjold og Værge!
Thi det er ikke noget Nyt, vi have for; ak, nei! det er en
gammel Strid, langvarig Rettergang fra Arilds Tid, fra
hiin Prinds Uffo med det rustne Sværd --- han greb til
Værget, da det gjaldt --- og saa igjennem Tiders, lange
Tiders Vexel ned indtil de sidste Offere, der faldt, rigt
smykkede med Kampens Roser. Det er en gammel Strid,
thi udaf Oltids Taager, Fortids Skumring samler sig en  % PRINTED AS IS: „Oltids“ for Oldtids (p. 15; crop J21)
Hær, en talrig, mægtig Hær, der strider for vor Ret; det
er en Forbundshær af stærke Minder: o, lad den være nær
paa Kampens Dag, i Farens Stund, den bryde frem med
samlet Styrke og med Kraft og Ild, som Flammer i en Luthers
Aand, som Sværdets Lyn i Gustav Adolphs Haand, at lyse
paa vor Vei til Maalet.

\end{document}
